Krátké zásady pro transparentní a odpovědné nasazení AI v předmětu, zaměřené na komunikaci se studenty a garanty:

\begin{itemize}
  \item Ve sylabu a na stránce předmětu jasně uveďte, které AI nástroje jsou povoleny/zakázány, jak mají studenti deklarovat použití AI a jak bude použití ovlivňovat hodnocení (viz šablony v kapitole 2.6 a příloha 9.3). % internal cross‑ref
  \item Požadujte jednoduchou provenance: u úloh, kde bylo AI použito, nechte studenty připojit krátké shrnutí co nástroj dělal, jaký prompt použili a které části výsledku upravili (general background knowledge).
  \item Zajistěte lidský dohled pro klíčové situace: kritické hodnocení výstupů AI provádí vyučující nebo asistent, zejména tam, kde výstupy mohou obsahovat chybné nebo citlivé informace.
  \item Logování a audit: pokud používáte enterprise nebo integrované řešení, aktivujte auditní logy a retenci provenance (dotazy/odpovědi, verze promptů/modelů) pro účely vyšetřování incidentů a zpětné kontroly \cite{kb_ai_101_tipu_pdf_04e4a115_page_3_1}.
  \item Informujte studenty o omezeních nástrojů (např. riziko „halucinací“) a o tom, jak ověřovat citace a fakta (general background knowledge). Pro technické integrace do testů a formulářů zvažte konkrétní procesy autorizace nástrojů (např. Copilot ve Forms) a popište je v instrukcích pro kurz \cite{kb_ai_101_tipu_pdf_04e4a115_page_36_1}.
  \item Dohoda o ochraně dat: při sběru dotazů/odpovědí nebo logů uveďte účel zpracování a dobu retention; pokud jde o citlivá data studentů, řešte to přes DPIA a DPA dle kapitoly 2.2. % internal cross‑ref
\end{itemize}

Praktický tip pro vyučující: připravte jednoznačnou větu do sylabu („Používání AI v tomto předmětu je povoleno za těchto podmínek: ...“), a při pilotu zaznamenávejte případy použití a incidenty pro kvartální revizi politiky.